Dit is een samenvatting en duiding van de meditaties met Sri Nisargadatta Maharaj, gebaseerd op zijn kernleren, vooral bekend uit het werk 'I Am That'. Deze lessen en dialogen vormen geen geloofssysteem, maar wijzen op een radicale, directe en compromisloze weg naar zelfrealisatie. Nisargadatta Maharaj, een eenvoudige beedishop-eigenaar uit Bombay, werd een gerespecteerde spirituele leraar, wiens inzichten de intellectuele complexiteit doorbraken en rechtstreeks tot de rauwe, onontkoombare waarheid van het bestaan doordrongen.

Zijn leringen gaan niet over het verkrijgen van iets nieuws, maar juist over het afwerpen van dat wat vals is en het herkennen van dat wat altijd al is. De kern van Maharaj’s onderricht kan samengevat worden in zijn centrale instructie: mediteer op het ‘ik ben’. Dit is geen mantra of intellectuele puzzel, maar een directe verwijzing naar het meest fundamentele en onmiskenbare aspect van het eigen bestaan: het ‘ik ben’.

Voor Maharaj is het ‘ik ben’ de fundamentele waarheid; het allereerste subtiele roeren van het bewustzijn, nog voor er een gedachte of gevoel is, of een identificatie als ‘ik ben een man’, ‘ik ben gelukkig’, of ‘ik zoek’. Er is enkel het onherleidbare besef van ‘ik ben’. Dit pure besef van aanwezigheid, van zijn, vormt de kern van onze ervaring. Het ‘ik ben’ is de eerste identificatie, de initiële manifestatie van het absolute, de ongemanifesteerde werkelijkheid. Het is de eerste stap naar manifestatie, het pure weten nog voordat het wordt ingevuld met inhoud. Maharaj benadrukt dat dit ‘ik ben’ de enige deur is naar zelfrealisatie: door zonder ophouden steeds opnieuw de aandacht terug te brengen naar dit rauwe besef van bestaan, begint men zich los te maken van alles wat erbij verzonnen of opgelegd is.

Het ‘ik ben’ is de stille, onveranderlijke getuige van alle gedachten, emoties, gevoelens en de buitenwereld. Het is het bewustzijn waarin alles verschijnt en weer verdwijnt, maar blijft zelf onaangeraakt door dat wat het waarneemt. De meditatie waarover Maharaj spreekt is geen formele zitpraktijk met bepaalde technieken, maar een voortdurend verblijven in dit besef van ‘ik ben’. Het vraagt om het te herkennen, ernaar terug te keren telkens als de geest afdwaalt en het tot het primaire focuspunt van de aandacht te maken. Maharaj zei: “Geef al je aandacht aan het gevoel van ‘ik ben’, ontdaan van alle concepten.”

Een kernpunt in zijn onderricht is het ontmantelen van ons diepgewortelde geloof in de realiteit van het afzonderlijke individu en een objectieve wereld. Lichaam en geest zijn tijdelijke manifestaties, instrumenten waarmee het ‘ik ben’ de wereld ervaart; ze zijn voortbrengselen van het bewustzijn, niet de bron. Ze zijn vergankelijk, voorwaardelijk en onderhevig aan verandering. Het ‘ik ben’ daarentegen is eeuwig en onveranderlijk. Maharaj stelt: je bent niet in de wereld; de wereld is in jou. Dit radicale inzicht houdt in dat de waargenomen wereld geen externe, onafhankelijke realiteit is, maar een manifestatie, een projectie binnen bewustzijn zelf, vergelijkbaar met een droom. In de droom lijkt de droomwereld echt, maar is volledig een product van de geest van de dromer.

Alle lijden ontstaat door valse identificatie met het lichaam-geest-complex en het geloof in een afgescheiden, eindig zelf. Wanneer men ten onrechte gelooft: “ik ben dit lichaam”, “ik ben deze geest”, dan worden de pijn, beperkingen en de uiteindelijke dood van het lichaam-geest-systeem een persoonlijk lijden. Maar voorbij deze identificatie is er geen afzonderlijk individu. Het gevoel van ‘ik’ dat zich met naam, vorm en persoonlijke geschiedenis identificeert, is een constructie, een tijdelijke bedekking van het pure ‘ik ben’.
